\documentclass[11pt]{article}
\usepackage[margin=25mm]{geometry}
\usepackage{fontspec}
\setmainfont{Latin Modern Roman}
\usepackage{amsmath,amssymb,mathtools}
\usepackage{unicode-math}
\setmathfont{Latin Modern Math}
\usepackage{xurl,hyperref}
\hypersetup{colorlinks=true,urlcolor=blue}
\setcounter{secnumdepth}{0}
\setlength{\parindent}{0pt}
\setlength{\parskip}{5pt}
\setlength{\emergencystretch}{3em}
\newcommand{\sref}[2]{\hyperlink{source:#1:#2}{[S#2]}}
\newcommand{\cataloguescope}[1]{\paragraph{Recorded target.}#1}
\begin{document}
\section{Lonely Runner Conjecture}\label{statement}

\subsection{Definitions and mathematical statement}

\paragraph{Shared notation.}$\mathbb N=\{1,2,\ldots\}$, and $\mathbb Z,\mathbb Q,\mathbb R,\mathbb C$ denote the integers, rationals, reals and complex numbers. Any other conventions are specified locally.


Represent a unit-circumference circular track by \(\mathbb R/\mathbb Z\).
For \(x\in\mathbb R\), let \(\|x\|=\min_{m\in\mathbb Z}|x-m|\);
the circular distance between \(x+\mathbb Z\) and \(y+\mathbb Z\)
is \(\|x-y\|\). Let \(k\ge2\) runners start together at \(0\),
with pairwise distinct constant real speeds \(v_1,\ldots,v_k\).
At time \(t\ge0\), runner \(i\) is at \(tv_i+\mathbb Z\).
Runner \(i\) is lonely at time \(t\) if its circular distance from
each other runner is at least \(1/k\).
\paragraph{Statement recorded in the supplied source.}
For every \(k\ge2\), every choice of pairwise distinct
\(v_1,\ldots,v_k\in\mathbb R\), and every
\(i\in\{1,\ldots,k\}\), does there exist a time \(t_i\ge0\) such that
\[
          \|t_i(v_i-v_j)\|\ge\frac1k
                 \qquad\text{for every }j\ne i?
\] \sref{statement}{2}
\subsection{Short English statement}
When any number of runners start together with different constant speeds, must each runner eventually be at circular distance at least the reciprocal of the number of runners from every other runner?
\subsection{Sources}
\begin{description}
\item[\hypertarget{source:statement:1}{S1}] ULAM.AI and the UnsolvedMath Contributors, \emph{UnsolvedMath: A Curated Collection of Open Mathematics Problems}, catalogue entry. CC BY 4.0. \url{https://huggingface.co/datasets/ulamai/UnsolvedMath}.
\item[\hypertarget{source:statement:2}{S2}] Touch Sungkawichai and Tanupat Trakulthongchai, \emph{Eleven, twelve, and thirteen lonely runners}, arXiv:2604.23906v2 (2026), Section 1, Conjecture 1.1 for distinct real speeds. \url{https://arxiv.org/abs/2604.23906}.
\end{description}
\label{end:statement}
\end{document}
